\section{Computable Functions}

\begin{definition}[Computable Function]
Let $D\subseteq\R$. A function $f\colon D\to\R$ is
\emph{computable} if an oracle algorithm $\A$ exists such that,
for every $x\in D$, every rational name
$\alpha\colon\Nzero\to\Q$ satisfying
\[
 |\alpha(i)-x|<2^{-i}\qquad(i\in\Nzero),
\]
and every $n\in\Nzero$, the algorithm with access to $\alpha$ and
input $n$ halts and returns $\A(\alpha, n)\in\Q$ with
\[
 |\A(\alpha, n)-f(x)|<2^{-n}.
\]
Only finitely many oracle queries are made in each computation.
The definition extends componentwise to maps between subsets of
$\R^k$ and $\R^m$.
\end{definition}

\begin{remark}
The name $\alpha$ need not itself be computable. The same oracle
algorithm must work for \emph{every} valid rational name, including
names of noncomputable reals. Unlike a binary-digit oracle, a rational
name need not settle dyadic boundary cases.
\end{remark}

If $f\colon D\to\R$ is computable, then
$f(D\cap\Rc)\subseteq\Rc$: a computable name for $x\in D\cap\Rc$
can be supplied to the oracle for $f$.

\begin{example}[Basic Computable Functions]
For $r\in\Rc$, the maps $x\mapsto x+r$ and $x\mapsto rx$ are
computable. So are $|x|$, $\exp x$, $\sin x$, $\sqrt{x}$ on
$[0, \infty)$, and $\log x$ on $(0, \infty)$.
In several variables, $+$, $-$, $\times$, $\min$, and $\max$ are
computable, as is division on $\{(x, y)\in\R^2\mid y\ne0\}$.
\end{example}

\begin{theorem}[Closure under Composition]
Let $f\colon D\to E$ and $g\colon E\to\R^m$ be computable,
where $D\subseteq\R^k$ and $E\subseteq\R^\ell$.
Then $g\circ f$ is computable.
\end{theorem}

\begin{corollary}[Computable Function Operations]
If $f, g\colon D\to\R$ are computable, then so are
\[
 |f|, \quad f+g, \quad fg, \quad \min(f, g), \quad\max(f, g),
 \quad\exp\circ f.
\]
Moreover, $f/g$ is computable on $\{x\in D\mid g(x)\ne0\}$.
\end{corollary}

\begin{theorem}[Computability Implies Continuity]
Every computable $f\colon D\to\R^m$, with $D\subseteq\R^k$,
is continuous on $D$ in the relative topology.
\end{theorem}

\begin{corollary}[Discontinuous Tests]
The sign function $\operatorname{sgn}\colon\R\to\{-1, 0, 1\}$
and the predicates $x<y$ and $x=y$, interpreted as
$\{0, 1\}$-valued functions on $\R^2$, are not computable
on their full domains.
\end{corollary}

\begin{remark}
Continuity is necessary but not sufficient for computability:
for $c\in\R\setminus\Rc$, the constant function $x\mapsto c$
is continuous but not computable. A discontinuity at a point
in the domain therefore prevents computability there.
\end{remark}

\begin{example}[Discontinuous Functions and Their Domains]
The step function $x\mapsto\indicator{\{x>0\}}$ is not computable on
$\R$, but is computable on $\R\setminus\{0\}$: refine the
rational approximation until its sign is certified. The same
principle makes strict comparison computable on $x\ne y$.

Alternatively, outside the single-valued framework above, one may
allow a set of admissible outputs at the discontinuity:
\[
 \operatorname{sgn}'(x)=
 \begin{cases}
  \{1\}, &x>0, \\
  \{0\}, &x<0, \\
  \{0, 1\}, &x=0.
 \end{cases}
\]
This is a multivalued specification, not a discontinuous
single-valued function made computable on its original domain.
\end{example}

\begin{theorem}[Computability of the Exponential Function]\label{thm:exp}
The function $\exp\colon\R\to\R$ is computable.
\end{theorem}
\begin{proof}
First let $x\in[-\tfrac12, \tfrac12]$ and let $\alpha$ be a
rational name. For prescribed $n\in\Nzero$, choose $m=d=n+3$ and put
\[
 \widehat{x}=\alpha(m), \qquad
 P_d(t)=\sum_{j=0}^{d}\frac{t^j}{j!}, \qquad
 \widehat{y}=P_d(\widehat{x}).
\]
This is a finite rational computation. Since
$|x-\widehat{x}|<2^{-m}$ and $|\widehat{x}|<1$, the mean value
theorem and $e<3$ imply
\[
 |e^x-e^{\widehat{x}}|<3\cdot2^{-m}.
\]
For $d\ge3$, Taylor's remainder is bounded by
\[
 |e^{\widehat{x}}-P_d(\widehat{x})|
 \le\sum_{j=d+1}^{\infty}\frac1{j!}
 \le\sum_{j=d+1}^{\infty}2^{-j}=2^{-d}.
\]
Consequently,
\[
 |e^x-\widehat{y}|<4\cdot2^{-(n+3)}<2^{-n}.
\]
For an arbitrary $x\in\R$, query $\alpha(0)$ and choose
$k\in\Nzero$ with $(|\alpha(0)|+1)/2^k\le\tfrac12$.
Then $|x/2^k|<\tfrac12$, and a rational name of $x/2^k$ is
obtained by scaling the oracle answers. Write
$u=e^{x/2^k}\in(0, 2)$ and $N=2^k$. For any rational $v$ with
$|u-v|<1$,
\[
 |u^N-v^N|\le N3^{N-1}|u-v|.
\]
Requesting a sufficiently accurate $v$ and computing the rational
$v^N$ therefore approximates $e^x$ to arbitrary precision.
\end{proof}

\begin{theorem}[Noncomputability of Strict Comparison]
The map
\[
 (x, y)\longmapsto
 \begin{cases}1, &x<y, \\0, &x\ge y\end{cases}
 \qquad (x, y\in\R)
\]
is not computable.
\end{theorem}
\begin{proof}
If the map were computable, requesting error less than $1/4$ and
comparing the rational answer with $1/2$ would decide its Boolean
output exactly. Consider the rational names
\[
 \alpha(i)=0, \qquad \beta(i)=2^{-(i+1)}
 \qquad(i\in\Nzero),
\]
both naming $0$. The decision algorithm returns $0$ after
only finitely many oracle queries. Let $N$ be at least every
index queried from $\beta$, and define
\[
 \beta'(i)=
 \begin{cases}
 2^{-(i+1)}, &i\le N, \\
 2^{-(N+1)}, &i>N.
 \end{cases}
\]
It is a valid name for $y'=2^{-(N+1)}>0$, because
$|\beta'(i)-y'|<2^{-i}$ for $i\le N$ and the error is zero
otherwise. The finite computation sees exactly the same answers
on $(\alpha, \beta')$ and $(\alpha, \beta)$, so it still returns
$0$, whereas $0<y'$ requires output $1$.
\end{proof}

\begin{theorem}[Comparison of Unequal Real Numbers]
Strict comparison is computable on
\[
 \{(x, y)\in\R^2\mid x\ne y\}.
\]
\end{theorem}

\begin{remark}
When $x\ne y$, sufficiently accurate rational approximations
must eventually certify their order. Without this a priori
promise, neither equality nor strict comparison is decidable
on all of $\R^2$: their indicator maps are discontinuous
along $x=y$.
\end{remark}
\Needspace{8\baselineskip}
